\subsection{Certifying fee selection, projection, and minimal allowance}\label{sec:semantic64}
An exact allocation proof and a correct claim about the chosen parameterization are different obligations. The tariff recurrence remains exact for any supplied standard fee after enumerating every relative exception. Its canonical certificate additionally claims that the standard is a most frequent original fee, with the smallest rational mode used to break ties. The independent checker now counts the original fees and checks this rule before checking the exception inventory, subset coverage, every Bellman state, and the returned original policy. Choosing a mode minimizes the exception count; choosing the smallest tied mode makes the implementation deterministic without changing its exact optimal value.

For a robust certificate the checker first binds the externally requested physical input. It sorts the original fees by value and catalog index, and, for the submitted allowance $d_{\rm allow}$, independently visits every consecutive window of length $N-d_{\rm allow}$. It evaluates the direct sum of absolute deviations from each lower median. This differs from the optimizer's prefix-sum calculation. It selects the minimum tuple (distortion, median, left position), and checks the supplied retained indices, median, complete projected charge vector, exception indices, and distortion. The inner exact tariff proof must use precisely this projected vector; its standard fee is independently checked against that vector's own modes. Thus $d_{\rm tariff}\leq d_{\rm proj}\leq d_{\rm allow}$, even when the two standard fees differ.

Write $D_d$ for the minimum total absolute distortion at allowance $d$. The checker verifies $D_{d_{\rm allow}}\leq\varepsilon$ and, when $d_{\rm allow}>0$, $D_{d_{\rm allow}-1}>\varepsilon$. Since $D_d$ is nonincreasing, the latter inequality excludes every smaller allowance; it is not necessary to trust the optimizer's binary-search path. At allowance zero no predecessor check is needed. Recomputing these two projections takes $O(N\log N+(d_{\rm allow}+1)N)$ rational operations. The certified minimum concerns the total absolute fee distortion, not the smallest final interval, realized regret, exact exception count after re-centering, or computation time. All original policy, one-sided error-budget, and final interval checks remain in force. The compatibility verifier delegates to the independent checker rather than retaining a weaker second verification path.

The semantic regression suite checks 288 small projections against exhaustive retained-subset enumeration and accepts 577 valid certificates, including equivalent rational encodings. Ten adversarial certificates carry false parameterization or projection claims while retaining valid objective intervals. They include a nonmodal standard, the wrong tied mode, a nonoptimal retained window, an incorrect median, and a nonminimal allowance. An isolated replay confirms that the archived checker accepts all ten; the current checker rejects all ten. This demonstrates the particular semantic closure, rather than counting arbitrary corrupted bytes as evidence. Checker import-closure tests exclude the optimization modules. These finite tests support implementation validation; the proofs in the article establish the mathematical guarantees.

The projection-quality table in the article is generated solely from frozen inputs, records, and certificates. It contains all twelve adversarial robust-tariff requests. Five original-fee regrets are identified exactly by closed rational intervals; two remain bounded rather than identified, and five requests have no certificate. Times and proof sizes are the original recorded quantities. Additional comparator checks determine value information only and do not alter a comparator's recorded deadline or checking outcome. The machine-readable table includes record and certificate hashes and keeps these two uses of evidence separate.
